\section{Related Work}
\label{sec:related_work}

We review five lines of work that bound the contribution: underwater simulation,
autonomous racing and reproducible robotics benchmarks, underwater perception,
multi-AUV cooperation, and learning-based marine control. This framing follows
the reproducibility and system-level validation perspective of recent benchmark
work~\cite{tani2020benchmark,araujo2023rasreview}. Table~\ref{tab:system_comparison}
summarizes the positioning.

\begin{table*}[t]
\centering
\caption{Task- and protocol-level properties of representative systems. The
comparison covers only properties relevant to the benchmark contribution of this
work; it does not rank simulator fidelity, sensing realism or throughput, and a
dash is not a deficiency.}
\label{tab:system_comparison}
\mratablestyle
\newcommand{\mrahead}[2]{\begin{tabular}[b]{@{}c@{}}#1\\#2\end{tabular}}%
\begin{tabular}{@{}llcccccc@{}}
\toprule
System & Domain &
\mrahead{Timed race}{task} &
\mrahead{Agent}{contract} &
\mrahead{Privileged-state}{restriction} &
\mrahead{Independent}{evaluator} &
Multi-vehicle &
\mrahead{Team}{score} \\
\midrule

UUV Simulator~\cite{manhas2016uuvsim}
& Underwater & -- & -- & -- & -- & \checkmark & -- \\

Stonefish~\cite{cieslak2019stonefish}
& Underwater & -- & -- & -- & -- & \checkmark & -- \\

DAVE~\cite{zhang2022dave}
& Underwater & -- & -- & -- & -- & \checkmark & -- \\

HoloOcean~\cite{potokar2024holoocean}
& Underwater & -- & -- & -- & -- & \checkmark & -- \\

UNav-Sim~\cite{amer2023unavsim}
& Underwater & -- & -- & -- & -- & -- & -- \\

URoBench~\cite{huang2024urobench}
& Underwater & -- & \checkmark & -- & -- & -- & -- \\

MarineGym~\cite{chu2025marinegym}
& Underwater & -- & \checkmark & -- & -- & -- & -- \\

F1TENTH~\cite{okelly2020f1tenth}
& Ground & \checkmark & \checkmark & -- & \checkmark & \checkmark & -- \\

CARLA~\cite{dosovitskiy2017carla} / Leaderboard~\cite{carlaLeaderboard2026}
& Ground & -- & \checkmark & \checkmark & \checkmark & \checkmark & -- \\

\midrule

\textbf{\mra{} (this work)}
& Underwater &
\checkmark &
\checkmark &
\checkmark &
\checkmark &
\checkmark &
\checkmark \\

\bottomrule
\end{tabular}

\mratablenote
Criteria. \emph{Timed race task}: a closed-course task over an ordered gate or
lap sequence whose official result is completion time. \emph{Agent contract}: a
documented participant-facing interface fixing the observations a scored
controller consumes and the command it returns. \emph{Privileged-state
restriction}: a protocol rule that withholds or forbids ground-truth simulator
state in scored runs. \emph{Independent evaluator}: the official result is
produced outside the submitted controller, which can neither observe nor
influence it; this denotes separation from the participant, not third-party
administration. \emph{Multi-vehicle}: several controllable vehicles in one
scenario, assessed at system level (simulator or benchmark) and excluding
parallel independent environment instances. \emph{Team score}: an aggregate
result over all vehicles of one entrant.
\checkmark: explicitly supported by the cited system; --: not defined as part of
a task or evaluation protocol in the cited system, which does not imply the
system is technically incapable of it.
\end{table*}


\subsection{Underwater Robotics Simulation}
\label{sec:rw_simulation}

Underwater simulators differ in physics engines, vehicle models and exposed
sensors. UUV Simulator provides Gazebo-based hydrodynamics for intervention and
multi-robot scenarios~\cite{manhas2016uuvsim}; Stonefish provides marine-specific
physics and underwater sensors through ROS~\cite{cieslak2019stonefish}; and DAVE
supports acoustic and optical sensing for autonomous underwater vehicles in a
Gazebo ecosystem~\cite{zhang2022dave}. HoloOcean uses Unreal Engine for rendered
underwater scenes and exposes sonar, RGB, IMU, DVL and depth sensing for one or
more agents~\cite{potokar2022holoocean,potokar2024holoocean}. UNav-Sim targets
underwater navigation and vision-based autonomy in a photorealistic Unreal-based
environment~\cite{amer2023unavsim}. Recent work also adds learning-oriented
Stonefish extensions~\cite{grimaldi2025stonefishml}, while MarineGym emphasizes
GPU-accelerated hydrodynamics and parallel training~\cite{chu2025marinegym}.

Benchmark surveys show that these systems make different trade-offs among
fidelity, throughput, sensors and task abstractions~\cite{aldhaheri2025simreview}.
URoBench makes that difference explicit by comparing underwater robotics
simulators through reproducible navigation tasks~\cite{huang2024urobench}. These
systems are complementary to \mra{}: they determine \emph{what happens} when a
vehicle acts in water, whereas \mra{} specifies \emph{how a race result is
judged}---which gates count, in which order, under what observation restriction,
adjudicated by whom, and recorded in what form. \mra{} adds this protocol above
the replaceable backend boundary.

\subsection{Autonomous Racing and Reproducible Robotics Benchmarks}
\label{sec:rw_racing}

Racing benchmarks expose end-to-end behaviour under repeatable and adversarial
constraints~\cite{betz2022autonomous}. F1TENTH standardizes a scaled car,
circuits and an evaluation environment for continuous control and reinforcement
learning~\cite{okelly2020f1tenth,evans2024f1tenth}; its ground-vehicle and planar-LiDAR assumptions
do not transfer directly underwater. At full scale, A2RL frames racing as
multi-agent interaction at the limits of vehicle handling~\cite{hoffmann2026a2rl,a2rl2024league}.
In aerial robotics, AlphaPilot established ordered gate sequences for fast
perception and control~\cite{foehn2020alphapilot,foehn2022alphapilot}, and Swift
demonstrated high-performance drone racing with deep reinforcement learning
against a simulated opponent~\cite{kaufmann2023swift}. The FeelHippo gate mission
illustrates a related perception-and-control challenge for underwater vehicles
~\cite{franchi2020feelhippo}. CARLA~\cite{dosovitskiy2017carla}, OpenAI Gym
~\cite{brockman2016openaigym} and Isaac Gym~\cite{makoviychuk2021isaacgym}
further demonstrate how documented interfaces can support systematic comparison.

\mra{} adopts ordered gates and a documented observation--action contract, but
underwater gate acquisition is constrained by short optical range, vehicle
damping, persistent flow and intermittent acoustic bearing. Gate association and
sequence adjudication therefore become benchmark responsibilities rather than
implementation details.

\subsection{Underwater Perception and Sensor-Based Navigation}
\label{sec:rw_perception}

Underwater autonomy is shaped by limited optical range, the absence of satellite
positioning and drift in inertial or Doppler estimates. These constraints motivate
vision-based localization and control with real vehicles~\cite{manzanilla2019uuvnavigation},
perception-aware navigation that preserves visual objectives~\cite{xanthidis2021aquavis},
and image-processing systems for a Gate Pass mission~\cite{harada2022gatepass}.
Tani et al. provide a recent visual--acoustic benchmark perspective for underwater
robot navigation~\cite{tani2026visualacoustic}. Simulators consequently model
optical and acoustic sensing explicitly~\cite{cieslak2019stonefish,zhang2022dave,potokar2024holoocean},
while underwater acoustic links remain noisy and low-rate~\cite{stojanovic2009underwater}.

\mra{} does not propose a new perception algorithm. Its narrower methodological
point is that, when several gates are visible, association between camera
detections and noisy acoustic cues remains in the participant autonomy stack
(Section~\ref{sec:perception}), rather than being resolved by the benchmark.

\subsection{Multi-AUV Cooperation and Acoustic Communication}
\label{sec:rw_fleet}

Fleet operation underwater is constrained by acoustic bandwidth, propagation
delay, multipath, attenuation and uncertain relative localization
~\cite{stojanovic2009underwater}. Formation-control surveys show how these
constraints shape achievable coordination and dependence on relative-state
estimation~\cite{yang2021survey,yan2023formation}. Recent marine coordination
studies likewise emphasize the coupling between communication, localization and
vehicle motion~\cite{martorell2024marinecoordination}.

Consistent with this literature, \mra{} assumes neither shared ground-truth state
nor high-rate communication. Fleet mode supports several vehicles as one team,
independent onboard autonomy and referee state, referee-side proximity checks
hidden from controllers, team-level scoring, and an acoustic-inspired channel
with range-dependent latency and loss (Section~\ref{sec:fleet}). This channel is
an experimental transport abstraction, not a complete independently calibrated
underwater channel model.

\subsection{Learning-Based Underwater Control and Simulation}
\label{sec:rw_learning}

Learning-based marine control has been enabled by simulators with Gym-style
interfaces and sufficient throughput. MarineGym is a direct recent reference:
it combines Isaac Sim, GPU hydrodynamics, multiple UUV models, domain
randomization and benchmark tasks for station keeping, tracking and docking
~\cite{chu2025marinegym}. The broader pattern---a documented observation--action
interface plus parallel simulation---also underlies comparable learned control
in ground and aerial robotics~\cite{brockman2016openaigym,makoviychuk2021isaacgym,kaufmann2023swift}.

MarineGym is primarily a training platform; \mra{} is an evaluation framework
for ordered-gate racing with an enforced onboard-only boundary, independent
referee, standardized penalties, multi-vehicle operation and team scoring. A
policy trained in a high-throughput platform can therefore be evaluated under
the \mra{} contract without changing either system. The learning result in
Section~\ref{sec:eval_rq5} is consequently a supporting demonstration of this
interface, not a claim that \mra{} replaces a training simulator.

\subsection{Positioning of \mra{}}
\label{sec:rw_positioning}

Table~\ref{tab:system_comparison} compares public system descriptions along
physics and sensing, ordered-gate racing, referee independence, observation
contract, multi-vehicle support and team scoring. The table is not a ranking;
the listed simulators exceed \mra{} on dimensions omitted here, including
hydrodynamic fidelity, sensor breadth and training throughput.

The claim is deliberately narrow: among the systems reviewed, \mra{} combines an
underwater ordered-gate racing protocol with structurally independent evaluation,
an observation contract that excludes privileged state from control, and
team-level evaluation. We do not claim to be the first underwater benchmark, the
first underwater multi-robot environment, or a better simulator than any system
in the table.
